\honest{The Wonder of Evolution}{Eliezer Yudkowsky}{2007}

The wonder of evolution, I say, is that it works at all. If every design needs a designer, who designed the designer? Evolution ends the regress by ``being stupid and inefficient and working anyway.''

When I describe evolution as slow, random and blind, someone asks whether that is not what the creationists say: that you cannot make a watch by shaking its parts in a box. My reply is that evolution does not shake parts in a box. Complex machinery evolves in small steps, or by putting old machinery to new uses; if birds died out, the descendants of gliding squirrels might learn to fly. Parts that are useful on their own can come to depend on one another, until the machine breaks if any one is removed. \nb{This answer to ``irreducible complexity'' is an old one. Hermann Muller described such ``interlocking'' in 1918, as an expected result of evolution.}

My example is DNA and the ribosome, each useless without the other. I say that ``in this case we do know how it happened,'' and I end the same paragraph with ``Almost certainly.'' RNA can store information and can also do chemistry, so it came first. \nb{The hedge is the accurate one. RNA before DNA and protein has broad support, but how an RNA system became one based on DNA and proteins is still an open problem.} RNA does both jobs badly, and that is the point.

The first replicator may have been RNA, I say, because RNA's ingredients ``would have arisen naturally'' on the early Earth. \nb{This was an open problem when the post was written. In 2009 the chemists who found a route to two of RNA's four nucleotides wrote that it was ``far from obvious'' how they could have formed.} Evolution does not explain the origin of life, I add, since the first replicator did not come from another replicator. It was a crude accident that only had to happen once.

Then I warn against praising evolution too highly. ``Science has a very exact idea of the capabilities of evolution,'' I say, and anyone who praises it more is being inaccurate, not loyal to science. I show the error through an imagined speaker.

I conclude that evolution is neither a designer nor a guide to conduct that humans should imitate. For human intelligence to imitate it as a designer would be like a modern bacterium imitating the first replicator. \nb{Engineers were already imitating evolution on computers, with success: antennas designed by evolutionary algorithms flew on a NASA mission in 2006.} I end with T. H. Huxley: the ethical progress of society depends on combating the cosmic process. \nb{Huxley's sentence is about ethics. It supports the second half of the conclusion, not the first.}
